\week{W5}{Trajectory generation}{5}
% W5 (booklet v3, condensed). Numbers verified by vcheck/v2_w05.py.
% plot panel: \wfiveP{xmax}{ymin}{ymax}{expression in \x}{label}
\newcommand{\wfiveP}[5]{\begin{tikzpicture}[x=1pt,y=1pt,every node/.style={font=\sffamily\fontsize{5.5}{6.5}\selectfont}]
 \pgfmathsetmacro{\wfsx}{64/(#1)}\pgfmathsetmacro{\wfsy}{40/((#3)-(#2))}
 \begin{scope}[x=\wfsx pt,y=\wfsy pt]
  \clip (0,#2) rectangle (#1,#3);
  \draw[g2,line width=0.3pt] (0,#2) rectangle (#1,#3);
  \draw[g3,line width=0.3pt] (0,0)--(#1,0);
  \draw[line width=0.8pt,samples=90,domain=0:#1] plot (\x,{#4});
 \end{scope}
 \draw[line width=0.3pt,g3] (0,{(#2)*\wfsy}) rectangle ({#1*\wfsx},{(#3)*\wfsy});
 \node[anchor=east] at (-1,{(#2)*\wfsy}) {\pgfmathprintnumber{#2}};
 \node[anchor=east] at (-1,{(#3)*\wfsy}) {\pgfmathprintnumber{#3}};
 \node[anchor=north] at ({#1*\wfsx},{(#2)*\wfsy-1}) {\pgfmathprintnumber{#1}\,s};
 \node[anchor=south west,font=\sffamily\fontsize{6}{7}\selectfont\bfseries,inner sep=1pt] at (0,{(#3)*\wfsy}) {#5};
\end{tikzpicture}}
\newcommand{\wfiveTri}[3]{\par\noindent\hspace*{4mm}#1\hfill #2\hfill #3\par}
% cubic 10->50 in 2 s
\def\wfCq{10+30*\x*\x-10*\x*\x*\x}
\def\wfCv{60*\x-30*\x*\x}
\def\wfCa{60-60*\x}
% LSPB 120 deg in 2 s, a=300 (tB=0.2254, v=67.62)
\def\wfLq{ifthenelse(\x<0.2254,150*\x*\x,ifthenelse(\x<1.7746,67.62*(\x-0.1127),120-150*(2-\x)*(2-\x)))}
\def\wfLv{ifthenelse(\x<0.2254,300*\x,ifthenelse(\x<1.7746,67.62,300*(2-\x)))}
\def\wfLa{ifthenelse(\x<0.2254,300,ifthenelse(\x<1.7746,0,-300))}

\begin{multicols}{3}

\hd{Joint vs Cartesian space \pg{6--14, 42--47}}
\premise{start/goal pose of $\{T\}$ (+ via points), $t_F$ \emph{or} speed/accel limits}{$q_i(t),\dot q_i,\ddot q_i$ for \emph{every} joint}
Trajectory = path + timing \pg{7}. Motors always get joint $q,\dot q,\ddot q$ \pg{68}. Teach mode \pg{9}: record and replay, no kinematics. Constant-acceleration formulas ($s=ut+\frac12at^2$, $v=u+at$, area under speed graph = distance): basics p.\,\pageref{p:BASICS}.
\par\begin{tabularx}{\linewidth}{@{}>{\bfseries\raggedright\arraybackslash}p{11mm}>{\raggedright\arraybackslash}X>{\raggedright\arraybackslash}X@{}}\toprule
 & Joint space \pg{12} & Task space \pg{43}\\\midrule
 Give & IK at both ends; a curve $q_i(t)$ per joint & path of $\{T\}$ in $\{S\}$ (line, ellipse), ${}^S\mathbf T_T(t)$\\
 Cost & cheap (IK twice) & IK at every sample, fit, resample\\
 Tool path & curved, uncontrolled \pg{46} & as asked\\
 Check & joint limits at ends & workspace, singularity, IK branch \emph{all along} \pg{47}\\\bottomrule
\end{tabularx}
\begin{fbx}[Linear interpolation \pg{14, 31}]
$q(t)=q_0+\dfrac{q_F-q_0}{t_F-t_0}(t-t_0)$, $\dot q$ = const. Velocity jumps $0\to\dot q$ at the ends (infinite acceleration), so use cubics/blends.
\end{fbx}

\hd{Cubic polynomial \pg{15--20}}
$N{+}1$ unknowns need $N{+}1$ conditions \pg{15}: position and velocity at both ends $\Rightarrow$ cubic.
\begin{fbx}[Cubic, $t_0=0$ \pg{16}]
\[q=a_0+a_1t+a_2t^2+a_3t^3,\ \ \dot q=a_1+2a_2t+3a_3t^2,\ \ \ddot q=2a_2+6a_3t\]
Conditions $q(0)=q_0$, $q(t_F)=q_F$, $\dot q(0)=\dot q_0$, $\dot q(t_F)=\dot q_F$. $a_k$ in unit/s$^k$.
\[\begin{bmatrix}1&0&0&0\\1&t_F&t_F^2&t_F^3\\0&1&0&0\\0&1&2t_F&3t_F^2\end{bmatrix}\!\begin{bmatrix}a_0\\a_1\\a_2\\a_3\end{bmatrix}=\begin{bmatrix}q_0\\q_F\\\dot q_0\\\dot q_F\end{bmatrix}\]
\[\mathbf M^{-1}=\begin{bmatrix}1&0&0&0\\0&0&1&0\\-\frac3{t_F^2}&\frac3{t_F^2}&-\frac2{t_F}&-\frac1{t_F}\\[1mm]\frac2{t_F^3}&-\frac2{t_F^3}&\frac1{t_F^2}&\frac1{t_F^2}\end{bmatrix}\]
\use{$t_0\ne0$: use local clock $\tau=t-t_0$, $T=t_F-t_0$ (via segments, any move not at $t=0$).}
\end{fbx}
\begin{deriv}[four coefficients; $\Delta q=q_F-q_0$]
$t=0$: $a_0=q_0$, $a_1=\dot q_0$. At $t_F$:
\[a_2t_F^2+a_3t_F^3=\Delta q-\dot q_0t_F,\quad 2a_2t_F+3a_3t_F^2=\dot q_F-\dot q_0\]
Eliminate $a_2$ ($=(\Delta q-\dot q_0t_F-a_3t_F^3)/t_F^2$): $a_3t_F^2=(\dot q_F+\dot q_0)-2\Delta q/t_F$.
\[\boxed{a_3=\frac{-2\Delta q+(\dot q_0+\dot q_F)t_F}{t_F^3}},\ \ \boxed{a_2=\frac{3\Delta q-(2\dot q_0+\dot q_F)t_F}{t_F^2}}\]
Rest to rest: $a_2=3\Delta q/t_F^2$, $a_3=-2\Delta q/t_F^3$; peak speed $1.5\Delta q/t_F$ at $t_F/2$; end accelerations $\pm6\Delta q/t_F^2$ (not zero). \chk
\end{deriv}
\begin{recipe}[cubic between two points]
\begin{enumerate}
\item List the 4 conditions with units; local clock, $T=t_F-t_0$.
\item $a_0=q_0$, $a_1=\dot q_0$; $a_2,a_3$ from the boxes.
\item Write $q,\dot q,\ddot q$ with numbers.
\item Check $q(T)=q_F$, $\dot q(T)=\dot q_F$.
\end{enumerate}
\end{recipe}
\begin{trap}
$a_3$ divides by $T^3$ (20$^\circ\to$80$^\circ$, 3\,s: $a_3=-2(60)/27=-4.44$, not $-13.3$). Do not mix degrees and radians.
\end{trap}
\begin{worked}[slide pp.\,17--19: $10^\circ\to50^\circ$ in 2\,s, rest to rest]
$\Delta q=40$: $a_2=3(40)/4=30$, $a_3=-2(40)/8=-10$.
$q=10+30t^2-10t^3$, $\dot q=60t-30t^2$, $\ddot q=60-60t$. $q(2)=10+120-80=50$ \checkmark, $\dot q(2)=0$ \checkmark; peak $\dot q(1)=30^\circ$/s. \chk\\
$\dot q_F=10$, $20^\circ\to80^\circ$, 3\,s: $a_2=\frac{180-30}{9}=16.667$, $a_3=\frac{-120+30}{27}=-3.333$ \chk. Several joints \pg{20}: one cubic each, same $T$ ($\Delta q=40,90,-20^\circ$, $T=2$: $a_2=30,67.5,-15$; $a_3=-10,-22.5,5$) \chk.
\end{worked}
\wfiveTri{\wfiveP{2}{10}{50}{\wfCq}{$q$ [$^\circ$]}}{\wfiveP{2}{0}{30}{\wfCv}{$\dot q$ [$^\circ$/s]}}{\wfiveP{2}{-60}{60}{\wfCa}{$\ddot q$ [$^\circ$/s$^2$]}}
\textbf{W13 sample Q9} (prismatic cubic, $0.02\to0.15$\,m in 0.8\,s) and \textbf{Q8} (via-point cubics): solved on p.\,\pageref{p:SAMPLE}.

\hd{Via points \pg{21--28}}
\premise{points $q_k$ at times $t_k$, end speeds (usually 0)}{cubic per segment, own clock $\tau\in[0,T_k]$}
Cubic between adjacent points; each via point needs a speed \pg{24}. Rest-to-rest at every point = stops everywhere \pg{25}. Options \pg{26}: (1) user gives speeds; (2) simple formula; (3) speeds chosen for \textbf{continuous acceleration}.
\begin{fbx}[Option 2: simple rule \pg{27}]
Chord slope $s_k=\dfrac{q_{k+1}-q_k}{t_{k+1}-t_k}$. At a via point between segments $k{-}1,k$:
\[\dot q_k=\begin{cases}0 & s_{k-1}s_k\le0\\ \tfrac12(s_{k-1}+s_k) & \text{same sign}\end{cases}\qquad \dot q=0\text{ at start, end}\]
Gives continuous $q,\dot q$; acceleration jumps.
\end{fbx}
\begin{worked}[notebook Q2: $q=10,50,20,40,45^\circ$ at $t=0,5,10,15,20$\,s]
$s=(8,-6,4,1)$; speeds $(0,0,0,2.5,0)$, $T_k=5$.
\begin{tabular}{@{}lrrrr@{}}\toprule seg.&$a_0$&$a_1$&$a_2$&$a_3$\\\midrule
$q_0\to q_A$&10&0&4.8&$-0.64$\\$q_A\to q_B$&50&0&$-3.6$&0.48\\$q_B\to q_C$&20&0&1.9&$-0.22$\\$q_C\to q_F$&40&2.5&$-0.4$&0.02\\\bottomrule\end{tabular}\\
Check $q_C\to q_F$ at 5\,s: $40+12.5-10+2.5=45$ \checkmark. $\ddot q$ jumps $+2.4,-3.4,+2.0$ at the joins. \chk
\end{worked}
\sh{Option 3: continuous acceleration \pg{28}, notebook Q3}
Slide: ``constant acceleration $a_{via}$''; read as \textbf{$\ddot\theta_1=\ddot\theta_2$ at the via point} (W13 Q8). Unknowns $a_0..a_3,b_0..b_3$ (8); $\theta_1=\sum a_kt^k$, $\theta_2=\sum b_kt^k$ (W13 p.\,18).
\begin{fbx}[8 conditions]
\begin{tabularx}{\linewidth}{@{}rXrX@{}}
1&$\theta_1(0)=\theta_0$&5&$\theta_2(T_2)=\theta_g$\\
2&$\dot\theta_1(0)=0$&6&$\dot\theta_2(T_2)=0$\\
3&$\theta_1(T_1)=\theta_v$&7&$\dot\theta_1(T_1)=\dot\theta_2(0)$\\
4&$\theta_2(0)=\theta_v$&8&$\ddot\theta_1(T_1)=\ddot\theta_2(0)$\\
\end{tabularx}
\end{fbx}
\begin{twoway}{A. Brute force}{B. Via-speed shortcut}
8 linear equations; 1, 2, 4 give $a_0=\theta_0$, $a_1=0$, $b_0=\theta_v$; solve 5 left.
\tcblower
Via speed $w$ unknown (7 holds); each segment an ordinary cubic with known end speeds; 8 fixes $w$.
\end{twoway}
\begin{deriv}[via speed $w$; $d_1=\theta_v-\theta_0$, $d_2=\theta_g-\theta_v$]
Seg.\,1 (end speeds $0,w$): $a_2=\frac{3d_1-wT_1}{T_1^2}$, $a_3=\frac{-2d_1+wT_1}{T_1^3}$. Seg.\,2 ($w,0$): $b_1=w$, $b_2=\frac{3d_2-2wT_2}{T_2^2}$, $b_3=\frac{-2d_2+wT_2}{T_2^3}$.
\[\ddot\theta_1(T_1)=2a_2+6a_3T_1=\tfrac{-6d_1+4wT_1}{T_1^2},\quad \ddot\theta_2(0)=2b_2=\tfrac{6d_2-4wT_2}{T_2^2}\]
Equate: $\boxed{w=\dfrac{3[T_1^2d_2+T_2^2d_1]}{2T_1T_2(T_1+T_2)}}$; $T_1=T_2$: $\boxed{w=\dfrac{3(\theta_g-\theta_0)}{4T}}$. Shared $\ddot\theta=3(\theta_0+\theta_g-2\theta_v)/T^2$. \chk
\end{deriv}
\begin{recipe}[two cubics through a via point]
\begin{enumerate}
\item Own clocks; list the 8 conditions.
\item $a_0=\theta_0$, $a_1=0$, $b_0=\theta_v$.
\item Cubics with end speeds $(0,w)$, $(w,0)$.
\item Equate accelerations at the via point; solve $w$.
\item Check positions, speeds, accelerations match.
\end{enumerate}
\end{recipe}
\begin{worked}[$\theta=10^\circ,40^\circ,20^\circ$, $T_1=T_2=1$\,s]
$w=3(10)/4=7.5$. Seg.\,1 $10+82.5t^2-52.5t^3$; seg.\,2 $40+7.5t-75t^2+47.5t^3$. $\theta_1(1)=40$, $\dot\theta_1(1)=165-157.5=7.5$, $\ddot\theta_1(1)=-150=2b_2$, $\theta_2(1)=20$, $\dot\theta_2(1)=0$ \checkmark. $T_2=2$: $w=25$; seg.\,1 $10+65t^2-35t^3$, seg.\,2 $40+25t-40t^2+11.25t^3$. \chk
\end{worked}

\hd{Quintic and time scaling \pg{29--30}}
\begin{fbx}[Quintic \pg{29}]
\[q=\textstyle\sum_{k=0}^5a_kt^k,\quad \dot q=\sum ka_kt^{k-1},\quad \ddot q=\sum k(k{-}1)a_kt^{k-2}\]
6 conditions: $q,\dot q,\ddot q$ at $0$ and $t_F$.
\[\begin{bmatrix}1&0&0&0&0&0\\1&t_F&t_F^2&t_F^3&t_F^4&t_F^5\\0&1&0&0&0&0\\0&1&2t_F&3t_F^2&4t_F^3&5t_F^4\\0&0&2&0&0&0\\0&0&2&6t_F&12t_F^2&20t_F^3\end{bmatrix}\!\begin{bmatrix}a_0\\a_1\\a_2\\a_3\\a_4\\a_5\end{bmatrix}=\begin{bmatrix}q_0\\q_F\\\dot q_0\\\dot q_F\\\ddot q_0\\\ddot q_F\end{bmatrix}\]
\use{zero acceleration at start and stop.}
\end{fbx}
\begin{deriv}[rest to rest, zero end acc. \added]
$a_0=q_0$, $a_1=a_2=0$. With $u=t_F$: $a_3u^3{+}a_4u^4{+}a_5u^5=\Delta$, $3a_3u^2{+}4a_4u^3{+}5a_5u^4=0$, $6a_3u{+}12a_4u^2{+}20a_5u^3=0$. Third: $a_3=-2a_4u-\frac{10}3a_5u^2$; second: $a_4=-\frac52a_5u$, $a_3=\frac53a_5u^2$; first: $a_5u^5/6=\Delta$.
\[\boxed{a_3=\tfrac{10\Delta}{t_F^3},\ \ a_4=-\tfrac{15\Delta}{t_F^4},\ \ a_5=\tfrac{6\Delta}{t_F^5}}\]
Peak speed $1.875\Delta/t_F$; $|\ddot q|_{max}=5.77\Delta/t_F^2$. $0\to15^\circ$, 1\,s: $(150,-225,90)$, $\dot q(1)=\ddot q(1)=0$ \checkmark. \chk
\end{deriv}
Time scaling \added: $q=q_0+\Delta q\,s(\tau)$, $\tau=(t-t_0)/T$; $s_{cubic}=3\tau^2-2\tau^3$, $s_{quintic}=10\tau^3-15\tau^4+6\tau^5$; same $s$ on all joints $\Rightarrow$ they finish together.
\par\begin{tabularx}{\linewidth}{@{}>{\raggedright\arraybackslash}Xcccc@{}}\toprule
 & lin. & cubic & quint. & LSPB$^*$\\\midrule
 peak $\dot q/(\Delta q/T)$ & 1 & 1.5 & 1.875 & 1.17\\
 end $\ddot q\cdot T^2/\Delta q$ & $\infty$ & $\pm6$ & 0 & $\pm8$\\\bottomrule
\end{tabularx}
{\footnotesize $^*a=2a_{min}$. Polynomials \pg{30}: speed is a hump, so a higher peak than the trapezoid for the same $\Delta q,T$.} \chk

\hd{Linear segment with parabolic blends (LSPB) \pg{31--37}}
\premise{$\Delta q=q_F-q_0$, $t_F$ ($t_0=0$), blend acceleration $a$}{blend time $t_B$, coast speed $v$, $q(t)$}
Speed up at $a$ for $t_B$, coast at $v$, slow down: speed graph = trapezoid \pg{32}.
\slidefig[0.9\linewidth]{w05-lec}{33}{19 119 19 62}{pp.\,33--34: $q,\dot q,\ddot q$ and the three position equations}
\begin{fbx}[Slide equations \pg{33--34}, $t_0=0$]
\[q_{B1}=q_0+\tfrac12at_B^2,\ \ q_{B2}=q_{B1}+v(t_F-2t_B),\ \ q_F=q_{B2}+vt_B-\tfrac12at_B^2,\ \ t_B=\tfrac va\]
$v$ coast speed; $a>0$ if $q_F>q_0$, $-a$ otherwise. Segments: $0\le t\le t_B$: $q=q_0+\frac12at^2$, $\dot q=at$; coast: $q=q_0+v(t-\frac12t_B)$, $\dot q=v$; last: $q=q_F-\frac12a(t_F-t)^2$, $\dot q=a(t_F-t)$.
\end{fbx}
\begin{deriv}[blend time and existence \pg{36}]
Area under trapezoid: $\Delta q=v(t_F-t_B)$ (sum of the 3 slide equations agrees). $v=at_B$:
\[\boxed{at_B^2-at_Ft_B+\Delta q=0}\Rightarrow t_B=\frac{t_F}2-\frac{\sqrt{a^2t_F^2-4a\Delta q}}{2a}\]
Real root needs $\boxed{a\ge4\Delta q/t_F^2}$. Blends must not overlap, $2t_B\le t_F$: take the \textbf{minus} root. Equality: $t_B=t_F/2$, triangle, $v=2\Delta q/t_F$. If $v$ is given: $t_B=t_F-\Delta q/v$, $a=v/t_B$, needs $\Delta q/t_F<v\le2\Delta q/t_F$. \chk
\end{deriv}
\begin{recipe}[LSPB, given $\Delta q,t_F,a$]
\begin{enumerate}
\item Existence: $a\ge4\Delta q/t_F^2$, else ``not possible''.
\item $t_B$ (minus root); $v=at_B$; coast time $t_F-2t_B$.
\item Check $v(t_F-t_B)=\Delta q$.
\end{enumerate}
\end{recipe}
\begin{worked}[slide p.\,37: $120^\circ$ in 2\,s, $a=300^\circ$/s$^2$]
$4(120)/4=120\le300$ \checkmark. $t_B=1\pm\frac{\sqrt{300^2\cdot4-4\cdot300\cdot120}}{600}=1\pm0.7746=0.2254$ (reject $1.7746$). $v=300(0.2254)=67.62^\circ$/s. Check $67.62(1.7746)=120$ \checkmark; $q_{B1}=7.62^\circ$. $a=120$: $t_B=1$ triangle; $a=100$: disc.\,$<0$; $v=70$ given: $t_B=0.2857$\,s, $a=245$. \chk
\end{worked}
\wfiveTri{\wfiveP{2}{0}{120}{\wfLq}{$q$ [$^\circ$]}}{\wfiveP{2}{0}{80}{\wfLv}{$\dot q$ [$^\circ$/s]}}{\wfiveP{2}{-300}{300}{\wfLa}{$\ddot q$ [$^\circ$/s$^2$]}}
\begin{fbx}[Minimum time, rest to rest, $|\dot q|\le v_{m}$, $|\ddot q|\le a_{m}$ \added]
$\Delta\ge v_m^2/a_m$ (trapezoid): $T_{min}=\dfrac{\Delta}{v_m}+\dfrac{v_m}{a_m}$. Else triangle: $T_{min}=2\sqrt{\Delta/a_m}$, $v_{pk}=\sqrt{a_m\Delta}$.
\use{$v_m=60$, $a_m=120$: $\Delta=90$: $1.5+0.5=2.0$\,s; $\Delta=20$: $0.816$\,s, peak $49.0$. Cubic needs $T\ge\max(1.5\Delta/v_m,\sqrt{6\Delta/a_m})=2.25$\,s for $90^\circ$. \chk}
\end{fbx}
\begin{trap}
Plus root ($t_B=1.7746$); $t_F$ for $t_F^2$ in $a_{min}$; ``bigger $a$ gives faster coast'' (it lowers $v$, toward $\Delta q/t_F$); sign of $a$ when $q_F<q_0$.
\end{trap}
\begin{slideerr}[pp.\,33--34]
General-$t_0$ form $q_{B1}=q_0+a(t_B-t_0)^2/2$ is only consistent for $t_0=0$ ($t_0=1,t_F=3$, $120^\circ$ gives roots $2.049,-0.049$, not $0.2254$). Use $T=t_F-t_0$ and a local clock.
\end{slideerr}

\hd{Several segments with blends \pg{38--41}}
Path passes \emph{near} the via points (rounded corners) \pg{40}. Given $\theta_k$, durations $t_{djk}$, $|\ddot\theta_j|$; find blend times $t_j$ (Craig \S7.3, \added).
\begin{fbx}[Blend times, $\theta_1\dots\theta_n$]
First: $t_1=t_{d12}-\sqrt{t_{d12}^2-\frac{2(\theta_2-\theta_1)}{\ddot\theta_1}}$, $v_{12}=\dfrac{\theta_2-\theta_1}{t_{d12}-\frac12t_1}$, $\ddot\theta_1=\operatorname{sgn}(\theta_2{-}\theta_1)|\ddot\theta_1|$.\\
Interior: $v_{jk}=\dfrac{\theta_k-\theta_j}{t_{djk}}$, $\ddot\theta_j=\operatorname{sgn}(v_{jk}-v_{ij})|\ddot\theta_j|$, $t_j=\dfrac{v_{jk}-v_{ij}}{\ddot\theta_j}$.\\
Last: $t_n=t_{d,n-1\,n}-\sqrt{t_{d,n-1\,n}^2+\frac{2(\theta_n-\theta_{n-1})}{\ddot\theta_n}}$, $v_{n-1\,n}=\dfrac{\theta_n-\theta_{n-1}}{t_{d,n-1\,n}-\frac12t_n}$, $\ddot\theta_n=\operatorname{sgn}(\theta_{n-1}{-}\theta_n)|\ddot\theta_n|$.\\
Linear parts: first $t_{12}=t_{d12}-t_1-\frac12t_2$; interior $t_{jk}=t_{djk}-\frac12t_j-\frac12t_k$; last $t_{n-1\,n}=t_{d,n-1\,n}-\frac12t_{n-1}-t_n$. Corner miss $|\ddot\theta_j|t_j^2/8$.
\end{fbx}
Example $0^\circ,40^\circ,10^\circ$, $t_d=2,2$, $|\ddot\theta|=50$: $t_1=0.4508$, $v_{12}=22.54$; $t_3=0.3267$, $v_{23}=-16.33$; $t_2=0.7775$; reaches $36.22^\circ$ (miss $3.78^\circ$). \chk

\hd{Cartesian paths, RR arm \pg{42--68}}
Use for a prescribed path or obstacles \pg{43}. Describe ${}^S\mathbf T_T(t)$; for IK ${}^B\mathbf T_T={}^B\mathbf T_S\,{}^S\mathbf T_T(t)$ \added. Do not interpolate rotation matrices entrywise \pg{65}.
\textbf{Joint-linear $\ne$ Cartesian-linear.} RR $L_1{=}4,L_2{=}3$, $(0,4.75)\to(4.75,0)$: joint-linear $q_1:51.07^\circ\to-38.93^\circ$, $q_2=95.83^\circ$ const; tool on an arc, midpoint $(3.359,3.359)$ vs line midpoint $(2.375,2.375)$, $1.39$\,m off. \chk
\begin{recipe}[check a Cartesian line first \pg{47}]
\begin{enumerate}
\item Workspace: $|L_1-L_2|\le r\le L_1+L_2$ on every point ($L_1{=}1,L_2{=}0.7$, line $y=0.2$: ends $r{=}1.02$ ok, middle $0.2<0.3$ hole). \chk
\item Singularity: $\det\mathbf J=L_1L_2\sin q_2$; joint rates $\propto1/\sin q_2$ ($L_1{=}L_2{=}1$, 1\,m/s on $y{=}0$: max rate $1.15,3.20,10.0,31.6$ rad/s at $x=1,1.9,1.99,1.999$). \chk
\item Same IK branch at both ends, else the arm jumps.
\end{enumerate}
\end{recipe}
\begin{recipe}[task-space trajectory \pg{48}]
\begin{enumerate}
\item Path as a function ($y=mx+c$).
\item Sample: $t_F/N$ steps; constant speed $v$: $x=x_0+\frac{vt}{\sqrt{1+m^2}}$, $y=y_0+\frac{vmt}{\sqrt{1+m^2}}$.
\item IK each sample, \textbf{same branch}.
\item Fit a cubic (quintic if needed \pg{60}) per joint.
\item Resample the fit; FK to see the tool path.
\end{enumerate}
\end{recipe}
\begin{fbx}[Planar RR FK/IK \pg{56}]
\[x=L_1c_1+L_2c_{12},\ \ y=L_1s_1+L_2s_{12},\ \ \cos q_2=\tfrac{x^2+y^2-L_1^2-L_2^2}{2L_1L_2}\]
\[q_2=\atantwo(\pm\sqrt{1-\cos^2q_2},\cos q_2)\]
\[q_1=\atantwo(y,x)-\atantwo(L_2\sin q_2,L_1+L_2\cos q_2)\]
Non-constant speed: $\dot{\mathbf q}=\mathbf J^{-1}\dot{\mathbf x}$.
\end{fbx}
\begin{worked}[slide pp.\,49--60: $L_1{=}4$, $L_2{=}3$, $(0,4.75)\to(4.75,0)$ at 3\,m/s, branch $q_2>0$]
Line $y=-x+4.75$, length $4.75\sqrt2=6.7175$, $T=2.2392$\,s; $x=3t/\sqrt2$, $y=4.75-x$; 10 sections $t_i=0.2239i$. IK table (computed, not transcribed):\\
{\scriptsize\begin{tabular}{@{}r@{\hspace{3pt}}r@{\hspace{3pt}}r@{\hspace{3pt}}r@{\hspace{3pt}}r@{\hspace{3pt}}r@{}}\toprule $i$&$t_i$&$x$&$y$&$q_1$ [$^\circ$]&$q_2$ [$^\circ$]\\\midrule
0&0&0&4.750&51.07&95.83\\1&0.224&0.475&4.275&41.48&105.71\\2&0.448&0.950&3.800&31.44&113.73\\3&0.672&1.425&3.325&20.75&119.76\\4&0.896&1.900&2.850&9.44&123.56\\5&1.120&2.375&2.375&$-2.13$&124.86\\6&1.344&2.850&1.900&$-13.18$&123.56\\7&1.567&3.325&1.425&$-22.85$&119.76\\8&1.791&3.800&0.950&$-30.48$&113.73\\9&2.015&4.275&0.475&$-35.84$&105.71\\10&2.239&4.750&0&$-38.93$&95.83\\\bottomrule\end{tabular}}\chk\\
$t_i=T/2$: $\cos q_2=\frac{2(2.375)^2-25}{24}=-0.5716$, $q_2=124.86^\circ$, $q_1=45^\circ-\atantwo(2.4616,2.2852)=-2.13^\circ$ (branch $-$: $92.13^\circ,-124.86^\circ$); FK: $(2.375,2.375)$ \checkmark. Cubic fits \pg{59}: $q_1=0.1346t^3-0.3375t^2-0.6182t+0.8858$, $q_2=0.0000t^3-0.4032t^2+0.9029t+1.6674$ (rad); fit error $0.46^\circ$, tool wanders $37$\,mm (quintic $7.7$\,mm). Other branch: $(128.93^\circ,-95.83^\circ)\to(38.93^\circ,-95.83^\circ)$, same path. Midpoint velocity $(3,-3)/\sqrt2$: $\det\mathbf J=12\sin q_2=9.846$, $\dot{\mathbf q}=(-0.893,0)$ rad/s. \chk
\end{worked}
\begin{slideerr}[p.\,59]
Slide prints $q_2=0.0004t^3-0.4043t^2+0.9036t+1.6674$; symmetry forces the cubic term to vanish (least squares: $0.0000t^3-0.4032t^2+0.9029t$). Method matters, not last digits.
\end{slideerr}

\hd{Exam-style questions}
Marks not published; $\sim$1 mark per step (conditions, formula, substitution, units, check).
\begin{examq}[LSPB \mk{marks not published}]
$75^\circ$ in $t_F=1.5$\,s, $a=200^\circ$/s$^2$. (a) Minimum $a$. (b) $t_B$, $v$. (c) Is $a=130$ possible?
\solution
(a) $a_{min}=4(75)/2.25=133.33$. (b) $\sqrt{90000-60000}=173.21$; $t_B=0.75-173.21/400=0.317$\,s; $v=at_B=\boxed{63.40^\circ/\text{s}}$; coast $0.866$\,s; check $63.40(1.183)=75.0$ \checkmark. (c) $130^2(2.25)-4(130)(75)=-975<0$: not possible. \chk
\end{examq}
\begin{examq}[Via points, simple rule \mk{marks not published}]
$10,30,70,80,60^\circ$ at $t=0,2,5,9,12$\,s. (a) Speed at 5\,s. (b) $a_2,a_3$ for $5\to9$\,s. (c) Acceleration jump at 5\,s.
\solution
$s=(10,13.333,2.5,-6.667)$. (a) $\frac12(13.333+2.5)=7.917^\circ$/s. (b) $T=4$, $\Delta=10$, $(\dot q_0,\dot q_F)=(7.917,0)$: $a_2=\frac{30-15.833(4)}{16}=-2.0833$, $a_3=\frac{-20+7.917(4)}{64}=0.18229$; $q(4)=80$ \checkmark. (c) Seg.\,$2\to5$ ($T=3$, $11.667\to7.917$): $a_2=2.9167$, $a_3=-0.7870$, end acc.\ $-8.333$; next start $2a_2=-4.167$; jump $\boxed{4.167^\circ/\text{s}^2}$. \chk
\end{examq}
\begin{examq}[RR tool on a circle \added\ \mk{marks not published}]
RR $L_1{=}0.6$, $L_2{=}0.4$\,m; tool once CCW round a circle $r{=}0.2$ centred $(0.5,0.3)$ in $T{=}4$\,s at constant speed, from $(0.7,0.3)$. (a) Reachable? (b) $q$ at $T/2$ ($q_2>0$). (c) Joint speeds at $t{=}0$.
\solution
$x=0.5+0.2\cos\frac{2\pi t}4$, $y=0.3+0.2\sin\frac{2\pi t}4$, speed $0.314$\,m/s. (a) Centre at $0.583$; points $0.383$ to $0.783$ from the base, inside $[0.2,1.0]$ \checkmark. (b) $(0.3,0.3)$: $\cos q_2=\frac{0.18-0.52}{0.48}=-0.7083$, $q_2=135.10^\circ$, $q_1=45^\circ-\atantwo(0.2824,0.3167)=\boxed{3.28^\circ}$. Others: $t{=}0$ $(-8.21^\circ,82.82^\circ)$; $t{=}1$ $(10.58^\circ,92.39^\circ)$; $t{=}3$ $(-29.95^\circ,122.80^\circ)$. (c) $\dot{\mathbf x}=(0,0.3142)$; $\mathbf J=\begin{bmatrix}-0.3000&-0.3857\\0.7000&0.1061\end{bmatrix}$, $\det=0.2381$; $\dot{\mathbf q}=\boxed{(0.509,-0.396)}$ rad/s. \chk
\end{examq}

\end{multicols}
